\section{Введение}

Ранее было показано, что задача о максимальном потоке решается за хоть сколько-то приемлимое время с помощью алгоритма Эдмондса-Карпа. 
Но глобально этот алгоритм особо ничем не отличается от алгоритма Форда-Фалкерсона, главная сложность там это показать, что мы имеем действительно асимптотику $\O\left(|V||E|^2\right)$.
Однако существует алгоритм, который нам позволит <<скинуть>> квадрат на множество вершин в худшем случае. Правда корректней тут говорить не сам алгоритм, а скорее схема, которая нам позволяет это сделать при должной имплементации. 
Эту схему называют \textit{схемой Диница}. Ключевой особенностью схемы является то, что мы ищем так называемые \textit{блокирующие потоки} в \textit{слоистой сети}, которую мы введем позднее. Центральным вопросом в этой схеме будет ответ на вопрос <<А как искать блокирующий поток?>>. 
На лекции будет приведен жадный способ их поиска, а также в качестве дополнительного материала будет рассмотрен способ удаляющего обхода, который как раз нам и позволит получить желаемую асимптотику $\O\left(|V|^2|E|\right)$.
Далее мы установим связь между <<пропускной способностью сети>> и количеством итераций схемы Диница.

Всюду далее мы считаем, что задана некоторая сеть $\mathcal{N} = (G, s, t, c)$.

\section{Блокирующие потоки}

\Def Пусть $f$ -- поток в сети $\mathcal{N}$. Ребро $e$ называется \textit{насыщенным}, если $f(e) = c(e)$.

\Def Поток $f$ в сети $\mathcal{N}$ будем называть \textit{блокирующим}, если его нельзя увеличить без введения обратных рёбер.

\Exe Покажите, что следующее определение блокирующего потока эквивалентно предыдущему. Пусть $f$ -- поток в сети $\mathcal{N}$. 
Если в $\mathcal{N}$ не существует пути от $s$ до $t$, который не проходил бы через насыщенные ребра потоком $f$, то поток $f$ называется \textit{блокирующим потоком}.

Читателю предлагается для себя ответить на вопрос: <<Какая связь между блокирующим потоком и максимальным? Является ли максимальный поток блокирующим? А наоборот?>>. В одну сторону очевидно: максимальный поток является блокирующим. Действительно, предположим, что это не так. 
Пусть $f$ максимальный, но не блокирующий поток. Тогда по определению блокирующего потока, существует путь не проходящий по насыщенным потоком $f$ рёбрам.
Возьмем такой путь $P$. Пустим вдоль него поток $f'$, тогда ясно, что $\forall e \in P \hookrightarrow f'(e) = c(P)$ и это поток в остаточной сети $\mathcal{N}_f$. 
Из \LmRef[леммы]{lm-flow-sum} следует, что $f + f'$ тоже есть поток и более того $|f + f'| > |f|$, что противоречит определению $f$. В обратную сторону существует контрпример:

\begin{figure*}[h]
    \centering
    \begin{tikzpicture}[>=Stealth, node distance=2cm, auto]
        
        \node[circle, draw, fill=white!20, minimum size=5mm] (A) at (0,0) {};
        \node[circle, draw, fill=white!20, minimum size=5mm] (B) at (2,1.5) {};
        \node[circle, draw, fill=white!20, minimum size=5mm] (C) at (2,-1.5) {};
        \node[circle, draw, fill=white!20, minimum size=5mm] (D) at (4,0) {};

        
        \draw[->] (A) -- (B) node[midway, above left] {1};
        \draw[->] (C) -- (D) node[midway, below right] {1};

        
        \draw[->, red] (A) -- (C) node[midway, below left] {1};
        \draw[->, red] (C) -- (B) node[midway, right] {1};
        \draw[->, red] (B) -- (D) node[midway, above right] {1};
    \end{tikzpicture}    
\end{figure*}

На рисунке показан блокирующий поток с величиной 1, однако, очевидно, если пусть поток по <<параллельным рёбрам>>, то он станет размера 4.

\section{Слоистая сеть}

\Def Слоистой сетью $(G_L\langle V_L, E_L \rangle, s, t, c)$ по сети $\mathcal{N}$ будем называть сеть построенную следующим образом:
\begin{equation*}
    V_d \overset{def}{=} \{ v: \delta(s, v) = d \}, d \in 0 \dots \delta(t)
\end{equation*}
\begin{equation*}
    V_L \overset{def}{=} \bigcup\limits_{d = 0}^{\delta(t)} V_d
\end{equation*}
$E_L$ получается же только лишь рёбрами, которые идут из $V_{d}$ в $V_{d + 1}$ для всякого $d$. Слоистую сеть порожденную сетью $\mathcal{N}$ будем обозначать $\mathcal{N}_L$ или $\mathcal{N}^L$.

Приведем пример слоистой сети. Ниже дан исходный граф, а далее его слоистая версия.

\begin{figure*}[h]
    \centering
    \makebox[\textwidth][c]{%
    \begin{tikzpicture}[
        node distance=1.5cm and 2cm,
        vertex/.style={circle, draw, thick, minimum size=12mm, inner sep=1pt, 
                    font=\large\bfseries\sffamily\boldmath, fill=white},
        edge/.style={-{Stealth[length=2.2mm, width=1.5mm]}, semithick, draw=black!70},
        layerbox/.style={draw=red!60!black, densely dashed, thick, rounded corners=5pt, 
                        fill=red!5, inner sep=8mm},
        layerlabel/.style={font=\small\bfseries, text=red!70!black},
        maintitle/.style={font=\Large\bfseries\sffamily, text=black!80},
        >=stealth,
        scale=0.52,
        transform shape,
        baseline=(current bounding box.center)
    ]

    \node[vertex] (s) at (0,0) {$s$};
    \node[vertex] (a) at (2,2) {$1$};
    \node[vertex] (b) at (2,-2) {$2$};
    \node[vertex] (c) at (4,3) {$3$};
    \node[vertex] (d) at (4,0) {$4$};
    \node[vertex] (e) at (4,-3) {$5$};
    \node[vertex] (t) at (6,0) {$t$};

    \draw[edge] (s) to (a);
    \draw[edge] (s) to (b);
    \draw[edge] (a) to (c);
    \draw[edge] (a) to (b);
    \draw[edge] (d) to (b);
    \draw[edge] (a) to (d);
    \draw[edge] (b) to (c);
    \draw[edge] (b) to (e);
    \draw[edge] (c) to (t);
    \draw[edge] (d) to (t);
    \draw[edge] (e) to (t);

    \end{tikzpicture}
    \hspace{0.6cm}%
    \begin{tikzpicture}[
        node distance=1.5cm and 2cm,
        vertex/.style={circle, draw, thick, minimum size=12mm, inner sep=1pt, 
                    font=\large\bfseries\sffamily\boldmath, fill=white},
        edge/.style={-{Stealth[length=2.2mm, width=1.5mm]}, semithick, draw=black!70},
        layerbox/.style={draw=red!60!black, densely dashed, semithick, rounded corners=3pt, 
                        fill=red!5, inner xsep=1.5mm, inner ysep=1mm},
        layerlabel/.style={font=\small\bfseries, text=red!70!black},
        maintitle/.style={font=\Large\bfseries\sffamily, text=black!80},
        >=stealth,
        scale=0.52,
        transform shape,
        baseline=(current bounding box.center)
    ]
        \node[vertex] (sL) at (0,0) {$s$};

        \node[vertex] (aL) at (3,1.5) {$1$};
        \node[vertex] (bL) at (3,-1.5) {$2$};

        \node[vertex] (cL) at (6,2.5) {$3$};
        \node[vertex] (dL) at (6,0) {$4$};
        \node[vertex] (eL) at (6,-2.5) {$5$};

        \node[vertex] (tL) at (9,0) {$t$};

        \draw[edge] (sL) to (aL);
        \draw[edge] (sL) to (bL);
        \draw[edge] (aL) to (cL);
        \draw[edge] (aL) to (dL);
        \draw[edge, dashed, gray] (aL) to (bL);
        \draw[edge, dashed, gray] (dL) to (bL);
        \draw[edge] (bL) to (cL);
        \draw[edge] (bL) to (eL);
        \draw[edge] (cL) to (tL);
        \draw[edge] (dL) to (tL);
        \draw[edge] (eL) to (tL);

        \begin{scope}[on background layer]
            \node[layerbox, fit=(sL)] (L0) {};
            \node[layerbox, fit=(aL)(bL)] (L1) {};
            \node[layerbox, fit=(cL)(dL)(eL)] (L2) {};
            \node[layerbox, fit=(tL)] (L3) {};
        \end{scope}
    \end{tikzpicture}
    }
\end{figure*}
Обратите внимание на пунктирные рёбра -- это те ребра, которые были в исходном графе, но не вошли в слоистую сеть.

Будем обозначать $\texttt{BUILD\_LAYERED(}\mathcal{N}\texttt{)}$ функцию, которая по сети $\mathcal{N}$ строит слоистую сеть.

\St[st-layered-shortest] Слоистая сеть содержит все кратчайшие пути от $s$ до $t$ и всякий путь в ней от $s$ до $t$ кратчайший.
\begin{proof}
    
    Рассмотрим произвольный путь $P = s \to v_1 \to v_2 \to \dots v_p \to t$ в слоистой сети. Заметим, что $p = \delta(t) - 1$, т.к. чтобы достичь $t$ мы должны пройти по всем слоям и последний обязательно будет под номером $\delta(t) - 1$.
    Значит вся длина пути $|P| = \delta(t)$, то есть он кратчайший.

    Покажем, что в оригинальной сети не было других кратчайших путей, которых не было бы в слоистой сети. Предположим, такой нашелся -- $P_0 = s \to v_1 \to v_2 \to \dots v_{\delta(t) - 1} \to t$.
    Вспомним, что всякий подпуть пути $P_0$ кратчайший. Тогда из этого следует, что $v_i \in V_i$.
    Более того, коль скоро $(v_i, v_{i + 1}) \in P_0$ и $v_i \in V_i, v_{i + 1} \in V_{i + 1} \Rightarrow$ по определению $E_L$, $(v_{i}, v_{i + 1}) \in E_L$. Аналогично, $(s, v_1)$ и $(v_{\delta(t) - 1}, t)$ лежат в $E_L$.
    Получается, все рёбра $P_0$ есть в слоистой сети, значит, $P_0$ находится в ней, что противоречит предположению.

\end{proof}

\section{Схема Диница}

Суть схемы достаточно проста. Строим слоистую по остаточной и находим в ней блокирующий поток.

\begin{algorithm}[H]
\caption{Схема Диница}
\begin{algorithmic}[1]

\Function{DINIC}{$\mathcal{N}$}
    \State $f \gets []$
    \While{$\text{есть блокирующий поток}$}
        \State $\mathcal{N}_L \gets \texttt{BUILD\_LAYERED(}\mathcal{N}_f\texttt{)}$
        \State $f' \gets \texttt{FIND\_BLOCK\_FLOW(}\mathcal{N}_L\texttt{)}$
        \State $f \gets f + f'$
    \EndWhile
    \State \Return $f$
\EndFunction
\end{algorithmic}
\end{algorithm}

\St \texttt{DINIC} возвращает максимальный поток.
\begin{proof}

    Алгоритм завершается, когда в остаточной пути нет блокирующего потока. Если нет блокирующего потока, значит в остаточной сети нет пути из $s$ в $t$, следовательно, по теореме Форда-Фалкерсона, $f$ -- максимальный.

\end{proof}

\Lm[lm-dinic-dist] После каждой итерации схемы Диница расстояние от истока к стоку в остаточной сети строго возрастает.
\begin{proof}

    Очевидно, что расстояние не может уменьшаться. Действительно, используя \LmRef[лемму]{lm-ek-dist} (если убрать из неё алгоритм Эдмондса-Карпа и работать в предположениях, что мы просто пускаем потоки вдоль кратчайших путей) получаем требуемое.

    Пусть на некоторый итерации поток был равен $f_0$ и после неё был найден блокирующий поток $f'$ в слоистой сети, порожденной остаточной сетью $G_f$.
    И пусть после этой итерации расстояние от $s$ до $t$ в $G_{f_0 + f'}$ не изменилось. Для определённости обозначим расстояние за $d$. И обозначим $G_0 \overset{def}{=} G_{f_0}, G_1 \overset{def}{=} G_{f_0 + f'}, \delta_0(s, v) \overset{def}{=} \delta_{G_{f_0}}(s, v), \delta_1(s, v) \overset{def}{=} \delta_{G_{f_0 + f'}}(s, v)$.

    Коль скоро расстояние не изменилось, существует путь $P$ длины $d$ в остаточной сети $G_1$. 

    Предположим, что он был в $G_0$. В силу того, что данный путь был кратчайшим от $s$ до $t$ в $G_0$, то, по \StRef[утверждению]{st-layered-shortest}, он находился в слоистой сети $G^L_0$.
    Но в силу того, что $P$ остался в сети $G_1$, ни одно ребро $P$ не насыщенно потоком $f_0$. Однако, по определению, $f_0$ блокирующий и всякий путь в $G_0^L$ (в том числе $P$), из $s$ в $t$ должен содержать рёбра насыщенные этим потоком. Противоречие.

    Утверждается, что существует некоторая вершина $v_0$, лежащая на данном пути, до которой расстояние увеличилось.
    Предположим противное $\Rightarrow$ расстояние до всех вершин на данном пути не изменилось. В $G_1$ есть только рёбра из $G_0$ или обратные к ним.
    Если бы $P$ состоял бы только из рёбер $G_0$, то мы приходим в прямое противоречие с предыдущим абзацем, т.к. получается что $P$ был в $G_0$. Значит $\exists (u, v) \in P$ обратное ребро к некоторому ребру из $G_0$.
    В силу того, что $P$ кратчайший, то $\delta_1(s, v) = \delta_1(s, u) + 1$.
    Более того, $(v, u)$ было в слоистой сети $G^L_0 \Rightarrow \delta_0(s, u) = \delta_0(s, v) + 1$ (т.к. слоистая сеть -- множество кратчайших путей, а подпуть кратчайшего кратчайший). Далее, комбинируя тот факт что расстояния не изменились:
    \begin{equation*}
        \begin{cases}
            \delta_1(s, v) = \delta_1(s, u) + 1 \\
            \delta_0(s, u) = \delta_0(s, v) + 1
        \end{cases}    
    \end{equation*}
    \begin{equation*}
            \delta_1(s, v) + \delta_0(s, u) = \delta_1(s, u) + \delta_0(s, v) + 1 + 1
    \end{equation*}
    \begin{equation*}
            \underbrace{\delta_1(s, v)}_{\delta_0(s, v)} + \delta_0(s, u) = \underbrace{\delta_1(s, u)}_{\delta_0(s, u)} + \delta_0(s, v) + 2
    \end{equation*}
    \begin{equation*}
            \delta_0(s, v) + \delta_0(s, u) = \delta_0(s, u) + \delta_0(s, v) + 2
    \end{equation*}
    \begin{equation*}
            0 = 2
    \end{equation*}
    Противоречие. Значит, существует вершина на пути $P$, до которой увеличилось расстояние. Возьмем самую последнюю из них на пути и обозначим $v_0$

    Предположим, что $v_0 \neq t \Rightarrow$ есть ребро $(v_0, u_0)$ на пути. Далее, $\delta_1(s, u_0) = \delta_1(s, v_0) + 1$, т.к. подпуть кратчайшего пути кратчайший.
    В силу выбора $v_0$, расстояние до $u_0$ не изменилось. Рассмотрим случаи:
    \begin{itemize}
        \item Ребро $(v_0, u_0)$ было в $G_0^L \Rightarrow \delta_0(s, u_0) = \delta_0(s, v_0) + 1$. 
        С другой стороны, $\delta_0(s, u_0) = \delta_1(s, u_0) \Rightarrow \delta_1(s, v_0) + 1 = \delta_0(s, v_0) + 1 \Rightarrow \delta_1(s, v_0) = \delta_0(s, v_0)$. Но ведь $\delta_1(s, v_0) > \delta_0(s, v_1)$. Противоречие.
        \item Ребро $(v_0, u_0)$ не было в $G_0^L$. В силу того, что $(v_0, u_0)$ появилось, то в $G_0^L$ было обратное к нему, то есть $(u_0, v_0)$. 
        Значит, $\delta_0(s, v_0) = \delta_0(s, u_0) + 1 = \delta_1(s, u_0) + 1 = \delta_1(s, v_0) + 2 \Rightarrow \delta_0(s, v_0) > \delta_1(s, v_0)$. Противоречие.
    \end{itemize}
    Значит, $v_0 = t \Rightarrow$ до $t$ расстояние увеличилось.

\end{proof}

\St Число итераций схемы Диница $\O(|V|)$.
\begin{proof}
    
    По \LmRef[лемме]{lm-dinic-dist}, расстояние от $s$ до $t$ строго возрастает. В тоже время максимальная длина пути от $s$ до $t$ не более $|V| - 1$. Получили требуемое.

\end{proof}

Теперь мы понимаем сколько итераций делает сама схема. 
Общее время работы определяется внутренним временем работы одной итерации цикла, которая состоит из построения слоистой сети и поика блокирующего потока и равна $\O(|E| + \tau)$, где $\tau$ -- время поиска блокирующего потока. Значит, общее время работы $\O(|V|(|E| + \tau))$.

Осталось понять как искать блокирующий поток. Вообще, вариантов как его искать много, однако мы ограничимся одним.

\subsection{Удаляющий обход}

Удаляющий обход заключается в том, что мы циклически запускаем DFS из $s$ в $t$ и пускаем поток вдоль пути, пока можем. 
Если в процессе DFS через ребро $e$ не смогли пустить, то мы и дальше не сможем через него пустить, значит, его можно не рассматривать далее и удалить.

Как это реализовать в компьютере? В списке смежности будем хранить указатель на первое неудаленное ребро. Если нам не удалось пустить поток через ребро, то двигаем указатель вперед.
Реализацию на языке C++ можно посмотреть \href{https://wiki.algocode.ru/index.php?title=Алгоритмы_поиска_блокирующего_потока}{здесь}. \textit{Настоятельно рекомендуется ознакомиться перед тем как читать далее}.

Сколько раз мы так сможем пустить DFS? В худшем случае у нас количество путей -- это количество рёбер, т.е. $\O(|E|)$ итераций.

Далее, сколько работает каждая итерация? Это вопрос чуть более нетривиальный и быть может не с первого раза понятный. Когда мы в DFS для очередной вершины рассматриваем ребро $e$, то мы либо пропускаем через него поток, либо понимаем, что не можем его пустить и двигаем указатель вперед. Рассмотрим подробнее ситуацию, когда мы не смогли пустить поток.
Как мы это понимаем? Выполняя рекурсивные шаги DFS, мы видим, что не можем достичь $t$, посредством того, что мы во всех путях рекурсии упираемся в другие вершины. А значит, как только мы вернемся по стеку рекурсии к ребру $e$, мы сдвинем для него указатель. 
Но важно, что перед этим мы еще сдвинем указатель для рёбер, которые были выше по стеку рекурсии. То есть для картинки ниже мы еще сдвинем указатели для $e_1, e_2, \dots e_n$ в процессе поднятия по стеку рекурсии к $e$.

\begin{figure*}[h]
    \centering
    \begin{tikzpicture}[
        node/.style={circle, draw, thick, minimum size=8mm, inner sep=0pt},
        dot/.style={inner sep=2pt, font=\large},
        arrow/.style={-{Stealth[length=3mm]}, thick},
        wavy/.style={decorate, decoration={snake, amplitude=1.2mm, segment length=4mm}},
        scale=0.75
    ]

    % Узлы
    \node[node] (s) {$s$};
    \node[node, right=1.8cm of s] (u) {$u$};
    \node[node, right=1.6cm of u] (v) {$v$};
    \node[node, right=1.6cm of v] (w1) {};
    \node[dot,   right=1.2cm of w1] (dots) {$\cdots$};
    \node[node,  right=1.2cm of dots] (wn) {};

    % Волнистая стрелка s -> u
    \draw[wavy, arrow] (s) -- (u);

    % Прямые рёбра с подписями
    \draw[arrow] (u)  -- node[above] {$e$}   (v);
    \draw[arrow] (v)  -- node[above] {$e_1$} (w1);
    \draw[arrow] (w1) -- node[above] {$e_2$} (dots.west);
    \draw[arrow] (dots.east) -- node[above] {$e_n$} (wn);

    \end{tikzpicture}
    
\end{figure*}

То есть если рассмотреть множество путей, которые мы рассмотрим, то будет примерно граф вида:

\begin{figure*}[h]
    \centering
    \begin{tikzpicture}[
        node/.style={circle, draw, thick, minimum size=8mm, inner sep=0pt},
        arrow/.style={-{Stealth[length=3mm, width=2mm]}, thick, shorten >=1pt, shorten <=1pt}
    ]

    \node[node] (s) {$s$};

    \node[node, above right=1.2cm and 2cm of s] (u1) {};
    \node[node, right=1.5cm of u1] (u2) {};

    \node[node, below right=1.2cm and 2cm of s] (v1) {};
    \node[node, right=1.5cm of v1] (v2) {};
    \node[node, right=1.5cm of v2] (v3) {};

    \node[node, right=2cm of s] (w1) {};
    \node[node, right=1.5cm of w1] (w2) {};
    \node[node, right=1.5cm of w2] (w3) {};

    \node[node, below right=2.8cm and 1.5cm of s] (x1) {};
    \node[node, right=1.5cm of x1] (x2) {};
    \node[node, right=1.5cm of x2] (x3) {};
    \node[node, right=1.5cm of x3] (t) {$t$};

    \draw[arrow, red] (s) -- (u1);
    \draw[arrow, red] (s) -- (v1);
    \draw[arrow, red] (s) -- (w1);
    \draw[arrow, blue] (s) -- (x1);

    \draw[arrow, red] (u1) -- (u2);
    \draw[arrow, red] (u1) -- (w2);
    \draw[arrow, red] (v1) -- (v2);
    \draw[arrow, red] (v2) -- (v3);

    \draw[arrow, red] (w1) -- (w2);
    \draw[arrow, red] (w2) -- (w3);

    \draw[arrow, blue] (x1) -- (x2);
    \draw[arrow, blue] (x2) -- (x3);
    \draw[arrow, blue] (x3) -- (t);
    \end{tikzpicture}
\end{figure*}

Где часть <<путей>>, где мы не смогли найти путь до $t$, и \textit{один простой путь от $s$ до $t$}. Соответственно, у тех путей, которые не ведут в $t$, мы удаляем ребра (помечены красным).
Время работы одной итерации оценивается количеством рёбер, которые мы посетили суммарно: количество <<удачных>> (синих) плюс <<неудачных>> (красных). Количество <<удачных>> рёбер  -- это длина простого пути от $s$ до $t$, т.е. $\O(|V|)$. Количество же <<неудачных>> есть ровно количество перемещений указателей. Для $i$-ой итерации мы обозначим количество сдвигов указателей $K_i$.

Далее, в силу того, что итераций в худшем случае есть $\O(|E|)$, то мы пройдем в худшем случае по $\O(|E|)$ простым путям. Каждый проход по простому пути, по рассуждениям выше, $\O(|V|)$. Значит <<суммарно по синим рёбрам в рамках всего алгоритма>> мы пройдем $\O(|V||E|)$ раз.
По красным же ребрам <<суммарно в рамках всего алгоритма>> мы пройдем $\O\left( \sum\limits_{i} K_i \right)$. То есть итоговое время работы алгоритма $\O\left( |V| |E| + \sum\limits_{i} K_i \right)$.
Однако в свою очередь, для каждого ребра мы в худшем случае один раз свдинем указатель, значит, $\sum\limits_{i} K_i \leq |E| \Rightarrow \O\left( \sum\limits_{i} K_i \right) = \O(|E|)$.
Значит итоговая асимптотика поиска блокирующего потока удаляющим обходом:
\begin{equation*}
    \O\left( |V| |E| + \sum\limits_{i} K_i \right) = \O\left( |V| |E| + |E| \right) = \O\left( |V| |E| \right)
\end{equation*}

\Cons Время работы схемы Диницы с использованием удаляющего обхода составляет $\O\left(|V|^2 |E|\right)$
\begin{proof}
    \begin{equation*}
        \O(|V|(|E| + \tau)) = \O\left(|V|^2|E|\right)
    \end{equation*}
\end{proof}

То есть мы получили асимптотически более быстрый алгоритм поиска максимального потока, чем алгоритм Эдмондса-Карпа, который работает за $\O\left( |V||E|^2 \right)$.

\section{Декомпозиция потока}

Следующая теорема показывает относительно интуитивные мысли о том, что поток <<распадается>> на пути и циклы.

\Def Пусть дано множество ребер $E$. Тогда
\begin{equation*}
    \chi_E (u, v) \eqd \begin{cases}
        1, (u, v) \in E \\
        -1, (v, u) \in E \\
        0, else
    \end{cases}
\end{equation*}

\Th{О декомпозиции} Пусть $\Ne$ -- сеть и $f$ -- поток в этой сети. Тогда найдутся простые пути из $s$ в $t$ $P_1, P_2, \dots, P_n$, а также циклы $C_1, C_2, \dots C_m$.
А также $\exists \alpha_1, \dots \alpha_n \in \N$ и $\beta_1, \dots, \beta_m \in \N$ такие, что:
\begin{equation*}
    f(u, v) = \sum\limits_{k = 1}^n \alpha_k \chi_{P_k}(u, v) + \sum\limits_{k = 1}^m \beta_k \chi_{C_k}(u, v)
\end{equation*}
То есть $f(u, v)$ есть <<сумма потоков>> путей, куда входит ребро $(u, v)$, а также быть может сумма <<потоков>> каких-то циклов.
Вообще говоря, <<поток>> в цикле называется циркуляцией (но этого знать необязательно) и это объект немного другой природы, который не рассматривается в нашем курсе.
\begin{proof}
    
Пусть из $s$ выходит хотя бы одно ребро $(s, v)$ с положительным потоком. Если такового нет, то $f \equiv 0$ и теорема тривиальна. Иначе перейдем по этому ребру в вершину $v$. 
Если $v = t$, то нашли путь $P$ из $s$ в $t$. Иначе, по закону сохранения, потока найдется ребро $(v, u)$ с положительным потоком. Если $u = t$, то нашли опять простой путь из $s$ в $t$. Иначе применим тот же алгоритм действий, что и для вершины $v$. И будем его применять, пока не достигнем, либо $t$, либо вершину, которую мы посетили на предыдущих шагах.

Рассмотрим случаи:
\begin{enumerate}
    \item Получили простой путь $P$ из $s$ в $t$. Положим $\alpha \eqd \min\limits_{e \in P} f(e)$. Очевидно, $\alpha > 0$. Более того, коль скоро емкости целочисленны $\Rightarrow \alpha \geq 1$. Уменьшим $f$ на $\alpha$. То есть новый рассматриваемый поток (и это действительно поток) на ребре $(u, v)$ будет равен $f(u, v) - \alpha, (u, v) \in P$ или $f(u, v) + \alpha, (v, u) \in P$ иначе $0$.
    \item Получили зациклившийся путь $P$. Вычленим из него цикл $C$. Положим $\beta \eqd \min\limits_{e \in C} f(e) \geq 1$. И аналогично прошлому пункту вычтем $\beta$ из $f$. Снова получили поток.
\end{enumerate}

Будем это делать до тех пор, пока $|f| > 0$. Далее, в следствие нашего алгоритма могли еще остаться циклы с $f(e) > 0$. Например, в такой ситуации:
\begin{figure}[H]
    \centering
    \begin{tikzpicture}[
        node style/.style={
            circle, 
            draw=black, 
            thick, 
            minimum size=8mm, 
            font=\bfseries\large
        },
        edge style/.style={
            ->, 
            thick, 
            >=Stealth, 
            shorten >=2pt, 
            shorten <=2pt
        }
    ]
        \node[node style] (s) at (0, 0) {$s$};
        \node[node style] (v1) at (3, 0) {$v_1$};
        \node[node style] (v2) at (6, 0) {$v_2$};
        \node[node style] (t) at (9, 0) {$t$};
        \node[node style] (v4) at (3, 3) {$v_4$};
        \node[node style] (v3) at (6, 3) {$v_3$};

        \draw[edge style] (s) -- node[above, font=\large] {$10/10$} (v1);
        \draw[edge style] (v1) -- node[below, font=\large] {$10/10$} (v2);
        \draw[edge style] (v2) -- node[below, font=\large] {$10/10$} (t);

        \draw[edge style] (v4) -- node[left, font=\large] {$5/5$} (v1);

        \draw[edge style] (v1) -- node[left, font=\large, pos=0.6] {$5/5$} (v3);

        \draw[edge style] (v3) -- node[above, font=\large] {$5/5$} (v4);

    \end{tikzpicture}

\end{figure}

Мы могли убрать путь $s \to v_1 \to v_2 \to t$ и тогда бы остался цикл $v_1 \to v_3 \to v_4 \to v_1$.

\begin{figure}[H]
    \centering
    \begin{tikzpicture}[
        node style/.style={
            circle, 
            draw=black, 
            thick, 
            minimum size=8mm, 
            font=\bfseries\large
        },
        edge style/.style={
            ->, 
            thick, 
            >=Stealth, 
            shorten >=2pt, 
            shorten <=2pt
        }
    ]
        \node[node style] (s) at (0, 0) {$s$};
        \node[node style] (v1) at (3, 0) {$v_1$};
        \node[node style] (v2) at (6, 0) {$v_2$};
        \node[node style] (t) at (9, 0) {$t$};
        \node[node style] (v4) at (3, 3) {$v_4$};
        \node[node style] (v3) at (6, 3) {$v_3$};

        \draw[edge style, dashed, gray] (s) -- node[above, font=\large] {$10/10$} (v1);
        \draw[edge style, dashed, gray] (v1) -- node[below, font=\large] {$10/10$} (v2);
        \draw[edge style, dashed, gray] (v2) -- node[below, font=\large] {$10/10$} (t);

        \draw[edge style, blue] (v4) -- node[left, font=\large] {$5/5$} (v1);

        \draw[edge style, blue] (v1) -- node[left, font=\large, pos=0.6] {$5/5$} (v3);

        \draw[edge style, blue] (v3) -- node[above, font=\large] {$5/5$} (v4);

    \end{tikzpicture}

\end{figure}

Найдем все такие циклы вычтем из $f$ аналогично пункту 2 выше и уже после всего этого получим, что $f \equiv 0$ (не величина, а прям функция потока). Тогда пусть мы получили пути $P_1, P_2, \dots P_n$ и с ними соответственно коэффициенты $\alpha_1, \alpha_2, \dots \alpha_n$ из пункта 1, а также циклы $C_1, C_2, \dots C_m$ и $\beta_1, \beta_2, \dots \beta_m$ из пункта 2 и предыдущего предложения.
Ясно, что в силу того, что мы в алгоритме итеративно вычитали из $f$ для соответствующих рёбер соответствующие коэффициенты $\{\alpha_i\}$ и $\{\beta_j\}$, то мы получили требуемые наборы.

\end{proof}

\Cons Пусть $\Ne$ -- сеть и $f$ -- поток в этой сети. Тогда найдутся простые пути из $s$ в $t$ $P_1, P_2, \dots, P_n$, а также циклы $C_1, C_2, \dots C_m$.
\begin{equation*}
    f(u, v) = \sum\limits_{k = 1}^n \chi_{P_k}(u, v) + \sum\limits_{k = 1}^m \chi_{C_k}(u, v)
\end{equation*}
Причем $n = |f|$. То есть иначе говоря существует декомпозиция по <<единичным путям и циклам>>.
\begin{proof}

    По теореме выше, найдутся простые пути из $s$ в $t$ $P_1, P_2, \dots, P_l$, а также циклы $C_1, C_2, \dots C_p$, а также $\exists \alpha_1, \dots \alpha_l \in \N$ и $\beta_1, \dots, \beta_p \in \N$, которые удовлетворяют условию теоремы.
    Теперь возьмем все такие $P_k$ и $C_k$ и проклонируем $\alpha_k$ и $\beta_k$ раз соответственно, но только уже с единичными коэффициентами. Получили требуемое.

    Покажем, что $n = |f|$. Заметим, что все простые пути начинают свой путь с рёбер вида $(s, v) \in E$. Далее:
    \begin{multline*}
        |f| \eqd \sum\limits_{(s, v) \in E} f(s, v) = \sum\limits_{(s, v) \in E} \left[ \sum\limits_{k = 1}^n \chi_{P_k}(s, v) + \sum\limits_{k = 1}^m \chi_{C_k}(s, v) \right] = \\
        = \sum\limits_{(s, v) \in E} \sum\limits_{k = 1}^n \chi_{P_k}(s, v) + \sum\limits_{(s, v) \in E}\sum\limits_{k = 1}^m \chi_{C_k}(s, v)
    \end{multline*}

    Покажем, что второе слагаемое есть ноль. Пусть $s \in C_k$, то есть $s$ находится в некотором цикле $\Rightarrow$ существует ребра $(p, s), (s, r) \in C_k$ -- входящее в $s$ и выходящее из него в цикле. Значит:
    \begin{equation*}
        \underbrace{\chi_{C_k} (p, s)}_{-1} + \underbrace{\chi_{C_k} (s, r)}_{1} = -1 + 1 = 0
    \end{equation*}
    Поэтому все слагаемые во втором слагаемом сократятся и получим:
    \begin{equation*}
        |f| = \sum\limits_{(s, v) \in E} \sum\limits_{k = 1}^n \chi_{P_k}(s, v)
    \end{equation*}
    Далее, если простой путь $P_k$ проходит по ребру $(s, v)$, то он точно не пройдет далее через некоторое другое ребро $(s, u)$, т.к. он бы пересек $s$ и не был бы простым. В то же время любой путь $P_k$ обязан проходить через какое-нибудь ребро $(s, v)$.
    Поэтому получается, что у нас пути $P_1, \dots P_n$ <<раскиданы>> по ребрам вида $(s, v)$, при этом один путь $P_k$ не может проходить через два таких ребра сразу $\Rightarrow$ двойная сумма выше есть просто количество рёбер, т.е. $n$:
    \begin{equation*}
        |f| = \sum\limits_{(s, v) \in E} \sum\limits_{k = 1}^n \chi_{P_k}(s, v) = n
    \end{equation*}
    \begin{equation*}
        |f| = n
    \end{equation*}
\end{proof}

\Exe Как нужно поменять условия теоремы о декомпозиции, чтобы она была верна, если образы функций потока и емкости вещественны? Верно ли тогда предыдущее следствие?

\section{Первая теорема Карзанова}
\begingroup

Если емкости вершин целочисленны (а в рамках нашей аксиоматики это так, хотя в общей постановке это важно), то на самом деле можно дать другую верхнюю оценку на количество итераций схемы Диницы.
В некоторых случаях данная оценка она может оценивать лучше, чем полученная в общем случае $\O(|V|)$. Однако данная оценка не выражается в наших привычных терминах $|E|$ и $|V|$, поэтому надо ввести некоторую дополнительную терминологию.

Всюду далее в этом разделе считаем, что зафиксирована некоторая транспортная сеть $\mathcal{N} = (G, s, t, c)$.

\newcommand{\p}{\mathbf{p}}

\renewcommand{\P}{\mathbf{P}}

\Def Пусть $\Ne$ -- сеть и $v \in V$. Тогда $c^+(v) \overset{def}{=} \sum\limits_{(u, v) \in E} c(u, v)$.

\Def Пусть $\Ne$ -- сеть и $v \in V$. Тогда $c^-(v) \overset{def}{=} \sum\limits_{(v, u) \in E} c(v, u)$.

\Def Пусть $\Ne$ -- сеть и $v \in V$. Тогда $\p(v) \overset{def}{=} \min\{c^+(v), c^-(v)\}$ -- потенциал вершины $v$.

О потенциале вершины можно думать как о некой верхней оценки потока, который может пройти через эту вершину.

\Def Пусть $\Ne$ -- сеть. Тогда потенциалом сети $\Ne$ будем называть
\begin{equation*}
    \P(\Ne) \eqd \sum\limits_{v \in V - s, t} \p(v)
\end{equation*}

На первый взгляд определение потенциала сети выше достаточно абстрактное и ни к чему необязывающее.
Ну действительно, с чего бы вдруг сумма каких-то оценок потоков в каждой вершины может что-то значить.
Однако певрая теорема Карзанова утверждает, что число итераций схемы Диница есть $\O\left(\sqrt {\P(\Ne)}\right)$. Перед доказательством данной теоремы мы докажем несколько утврерждений.

\Def $V_d \eqd \{ v: \delta(s, v) \}$

Покажем очевидные утверждения.

\St $\forall k \geq 2 \, \nexists (u, v) \in E: u \in V_d, v \in V_{d+k}$
\begin{proof}
    
    Действительно, если бы такое ребро существовало, то $v$ было бы на расстоянии $d + 1 < d + k \Rightarrow v \in V_{d + 1}$, противоречие.

\end{proof}

\St Верно следующее:
\begin{equation*}
    \sum\limits_{d = 1}^{\delta(s, t) - 1} \P(V_d) \leq \P(\Ne) 
\end{equation*}
\begin{proof}
    \begin{equation*}
        \sum\limits_{d = 1}^{\delta(s, t) - 1} \P(V_d) = \sum\limits_{d = 1}^{\delta(s, t) - 1} \sum\limits_{v \in V_d} \p(v) = \left\{ \bigcup_{d = 1}^{\delta(s, t) - 1} V_d \subset V \right\} = \sum\limits_{v \in V - s, t} \p(v) \leq \P(\Ne)
    \end{equation*}
\end{proof}

\Def Пусть $V' \subset V$. Тогда:
\begin{equation*}
    \P(V') \eqd \sum\limits_{v \in V' - s, t} \p(v)
\end{equation*}

\St Пусть $f_0$ -- максимальный поток в сети. Тогда $\forall d \in 1 \dots \delta(t) - 1 \hookrightarrow |f_0| \leq \P(V_d)$
\begin{proof}
    
    По следствию теоремы о декомпозиции, найдутся <<единичные>> простые пути из $s$ в $t$ $P_1, P_2, \dots, P_n$, а также циклы $C_1, C_2, \dots C_m:$
    \begin{equation*}
        f_0(u, v) = \sum\limits_{k = 1}^n \chi_{P_k}(u, v) + \sum\limits_{k = 1}^m \chi_{C_k}(u, v)
    \end{equation*}
    Далее, в силу того, что $V_d$ -- множество вершин на расстоянии $d \leq \delta(s, t)$, то всякий простой путь $P_k$ содержит хотя бы одну вершину из $V_d$. Значит, он уменьшает <<входную и выходную>> емкости хотя бы на 1. Значит все пути в сумме уменьшают <<входные и выходные>> емкости хотя бы на $n$.
    Из этого моментально следует, что $\P(V_d) - n \geq 0 \iff \P(V_d) \geq n$, иначе бы не хватило емкости. С другой стороны, по тому же следствию из теоремы о декомпозиции, $|f_0| = n \Rightarrow \P(V_d) \geq n = |f_0|$. Получили требуемое.

\end{proof}

\Lm Пусть $f_0$ -- максимальный поток в сети $\Ne$. Тогда верно следующее соотношение:
\begin{equation*}
    \delta_{\Ne}(s, t) \leq \frac{\P(\Ne)}{|f_0|} + 1
\end{equation*}
\begin{proof}
    
    По лемме выше, $\forall d \in 1 \dots \delta(s, t) - 1 \hookrightarrow |f_0| \leq \P(V_d)$. Тогда:
    \begin{equation*}
        (\delta_\Ne(s, t) - 1) |f_0| \leq \sum\limits_{d = 1}^{\delta_\Ne(s, t) - 1} \P (V_d) \leq \P(\Ne)
    \end{equation*}
    \begin{equation*}
        \delta_\Ne(s, t) - 1 \leq \frac{\P(\Ne)}{|f_0|}
    \end{equation*}
    \begin{equation*}
        \delta_\Ne(s, t) \leq \frac{\P(\Ne)}{|f_0|} + 1
    \end{equation*}

\end{proof}

\Lm Пусть $f$ -- поток в сети $\Ne$. Тогда $\P(\Ne) = \P\left(\Ne_f\right)$.
\begin{proof}

    Покажем, что $c^+$ и $c^-$ не изменились в остаточной сети. Покажем для $c^+$, т.к. для $c^-$ аналогично.

    Воспользуемся следствием теоремы о декомпозиции потока для $f$.
    Существуют <<единичные>> простые пути из $s$ в $t$ $P_1, P_2, \dots, P_n$, а также циклы $C_1, C_2, \dots C_m:$
    \begin{equation*}
        f(u, v) = \sum\limits_{k = 1}^n \chi_{P_k}(u, v) + \sum\limits_{k = 1}^m \chi_{C_k}(u, v)
    \end{equation*}

    Рассмотрим вершину $v \in V - s, t$ и рассмотрим все пути и циклы, полученные выше, которые проходят через неё.
    Рассмотрим $P_k$ или $C_k$ проходящий через $v$. Они обязаны пускать поток (или в случае цикла циркуляцию) следующим образом:

    \begin{figure}[H]
        \center
        \begin{tikzpicture}[
            node/.style={circle, draw, thick, minimum size=8mm, font=\large},
            edge label/.style={font=\small, inner sep=2pt}
            ]
                \node[node] (u) at (0,0) {$u$};
                \node[node, blue] (v) at (3,0) {$v$};
                \node[node] (w) at (6,0) {$w$};

                \draw[-{Stealth[length=3mm]}, thick] 
                    ([yshift=4pt] u.east) -- ([yshift=4pt] v.west)
                    node[midway, above, edge label] {$+1$};
                \draw[-{Stealth[length=3mm]}, thick] 
                    ([yshift=-4pt] v.west) -- ([yshift=-4pt] u.east)
                    node[midway, below, edge label] {$-1$};

                \draw[-{Stealth[length=3mm]}, thick] 
                    ([yshift=4pt] v.east) -- ([yshift=4pt] w.west)
                    node[midway, above, edge label] {$+1$};
                \draw[-{Stealth[length=3mm]}, thick] 
                    ([yshift=-4pt] w.west) -- ([yshift=-4pt] v.east)
                    node[midway, below, edge label] {$-1$};
            \end{tikzpicture}    
    \end{figure}
    То есть поток (циркуляция) еще обратные ребра создает. Получается, что после пропускания пути (цикла) через $v$ в неё \textit{входят} ребра, которые увеличивают входную емкость на $+1$ и $-1$. В сумме получаем $+1 - 1 = 0$, то есть <<единичный>> путь (цикл) проходящий через $v$ не меняет входную емкость. Значит, любая сумма их тоже не меняет $\Rightarrow$ и сам поток тоже ничего не меняет.

\end{proof}

\Th{Первая теорема Карзанова} Количество итераций схемы Диница равно $\O\left(\sqrt {\P(\Ne)}\right)$.
\begin{proof}
    
    Запустим схему Диница на $\sqrt {\P(\Ne)}$. Если больше запускать ее не надо, то это тривиальный случай. Допустим, еще можно запустить итерации. 
    
    Пусть после $\sqrt {\P(\Ne)}$ итераций мы получили поток $f$. По лемме выше, $\P(\Ne) = \P(\Ne_f)$. Пусть $f'$ -- максимальный поток в $\Ne_f$.
    
    В силу того, что после каждой итерации схемы расстояние от истока до стока увеличивается хотя бы на 1, то $\delta_{\Ne_f}(s, t) \geq \sqrt {\P(\Ne)}$.

    Далее, применим вышедоказанную лемму к сети $\Ne_f$ (вспомним, что в ней максимальный поток $f'$):
    \begin{equation*}
        \delta_{\Ne_f}(s, t) \leq \frac{\P\left(\Ne_f\right)}{|f'|} + 1
    \end{equation*}
    Далее, вспомним, что $\P(\Ne) = \P\left( \Ne_f \right)$:
    \begin{equation*}
        \delta_{\Ne_f}(s, t) \leq \frac{\P\left(\Ne\right)}{|f'|} + 1
    \end{equation*}
    С другой стороны мы оценили $\delta_{\Ne_f}(s, t)$ снизу, а значит:
    \begin{equation*}
        \sqrt {\P(\Ne)} \leq \delta_{\Ne_f}(s, t) \leq \frac{\P\left(\Ne\right)}{|f'|} + 1
    \end{equation*}
    \begin{equation*}
        \sqrt {\P(\Ne)} \leq \frac{\P\left(\Ne\right)}{|f'|} + 1
    \end{equation*}
    Теперь выразим оценку на $|f'|$:
    \begin{equation*}
        \frac{\P\left(\Ne\right)}{|f'|} + 1 \geq \sqrt {\P(\Ne)}
    \end{equation*}
    \begin{equation*}
        \frac{\P\left(\Ne\right)}{|f'|} \geq \sqrt {\P(\Ne)} - 1
    \end{equation*}
    \begin{equation*}
        \frac{|f'|}{\P\left(\Ne\right)} \leq \frac{1}{\sqrt {\P(\Ne)} - 1}
    \end{equation*}
    \begin{equation*}
        |f'| \leq \frac{\P\left(\Ne\right)}{\sqrt {\P(\Ne)} - 1} = \O\left(\sqrt {\P(\Ne)}\right)
    \end{equation*}
    \begin{equation*}
        |f'| = \O\left(\sqrt {\P(\Ne)}\right)
    \end{equation*}
    Далее, в силу того, что поток в сети $\Ne_f$ с каждой итерацией увеличивается хотя бы на 1, то в худшем случае нужно будет сделать еще $|f'| = \O\left(\sqrt {\P(\Ne)}\right)$ итераций.
    Значит суммарно мы сделаем $\sqrt {\P(\Ne)} + \O\left(\sqrt {\P(\Ne)}\right) = \O\left(\sqrt {\P(\Ne)}\right)$ итераций.

\end{proof}

\Cons Схема Диница с удаляющим обходом работает за $\O\left( |V||E| \sqrt{\P(\Ne)} \right)$.

\section{Алгоритм Хопкрофта-Карпа}

Наглядной демонстрацией применения первой теоремы Карзанова -- это алгоритм Хопкрофта-Карпа. Этот алгоритм решает задачу нахождения максимального паросочетания.
По своей сути он просто тривиален, но без знания теоремы Карзанова мы получими закономерный вопрос о состоятельности и нужности этого алгоритма. Однако давайте же его рассмотрим.

Пусть дан на вход двудольный граф $G\langle L \oplus R, E \rangle$. Далее, $V = L \cup R$. Построим на его основе сеть. Введем исток $s$ и соединим со всеми рёбрами.
Для каждого такого ребра $(s, u)$ положим $c(s, u) \eqd 1$. Далее, ориентируем ребра паросочетания так, чтобы они исходили из $L$ и входили в $R$ и для каждого такого ребра $(l, r)$ положим $c(l, r) = 1$.
Введём сток $t$ и соеденим вершины из $R$ с $t$ рёбрами емкостью $1$. Обозначим полученную сеть за $\Ne_0$. Сразу заметим, что она единичная. Очевидно, что множество рёбер, насыщенные максимальным потоком, представляют собой максильное паросочетание.
Также особое внимание уделим тому, что в этой сети $|E| + 2 |V|$ рёбер (исходные рёбра графа и еще ребра на соединение со стоком и истоком) и $|V| + 2$ вершины

\Exe В принципе корректность данного алгоритма более чем очевидна, однако попробуйте ее показать с помощью теоремы Бержа.

Если искать максимальный поток в такой сети с помощью схемы Диница с удаляющим обходом, мы, на первый взгляд, получаем асимптотику $\O\left(|V|^2(|E| + |V|)\right)$. Это, казалось бы, асимптотически медленнее, чем алгоритм Куна, работающий за $\O(|V|(|V| + |E|))$.

Однако, давайте посчитаем потенциал данной сети.

\begin{equation*}
    \P(\Ne_0) = \sum\limits_{v \in V - s, t} \p(v) = \sum\limits_{v \in V - s, t} \underbrace{\min \{ c^-(v), c^+(v) \}}_{1} = \Theta(|V|)
\end{equation*}

Далее, рассмотрим как работает удаляющий обход. Заметим, что после каждой итерации удаляющего обхода все ребра насыщенные блокирующим потоком забываются (в силу единичности сети). Значит, удаляющий обход работает ровно за $\O(|V| + |E|)$.

Получаем:
\begin{equation*}
    \O\left( (|V| + |E|) \sqrt{\P(\Ne_0)} \right) = \O\left( (|V| + |E|) \sqrt{|V|} \right)
\end{equation*}
Если граф <<адекватный>> (то есть $|V| = \O(|E|)$), то получаем еще более красивую асимптотику $\O\left(|E| \sqrt{|V|}\right)$.

\endgroup
